
\section{Selective Harmonic Response}

Let the complement correlation length scale with the parent radius,
\[
L_c(a)=aR_H\,\ell,
\]
where \(\ell\) is dimensionless. Since
\[
p_n=\frac{n(n+2)}{a^2R_H^2},
\]
we have
\[
L_c^2p_n=\ell^2n(n+2).
\]

In the effectively massless topographic sector, define
\[
K_n=p_nC_n
\]
with
\[
\boxed{
C_n(a)
=
c_s^2(a)
-
\chi(a)\left[1-\ell^2n(n+2)\right].
}
\]

If
\[
\boxed{
\frac18<\ell^2<\frac13,
}
\]
then
\[
1-3\ell^2>0
\]
for \(n=1\), whereas
\[
1-\ell^2n(n+2)<0
\]
for every \(n\ge2\). Hence increasing \(\chi\) softens the \(n=1\) response but stiffens all higher \(S^3\) scalar levels:
\[
C_1<c_s^2,\qquad
C_n>c_s^2\quad(n\ge2).
\]

This defines a selective positive-kernel branch within the reduced spatial ansatz. If
\[
0<C_1\ll1,
\]
then a sourced mode
\[
T_1\simeq\frac{S_1}{p_1C_1}
\]
can be strongly enhanced while the \(n\ge2\) responses remain regular. This positivity statement is narrower than a full perturbative stability result because the complete closed-background kinetic/gradient system has not yet been diagonalized.

For the next \(S^3\) scalar level,
\[
\frac{T_2}{T_1}
=
\frac{S_2}{S_1}
\frac{3C_1}{8C_2},
\]
so the response per unit source satisfies
\[
\left|
\frac{T_2/S_2}{T_1/S_1}
\right|
=
\frac{3C_1}{8C_2}
<
\frac38
\]
on the positive selective-filter branch, with much stronger suppression as \(C_1/C_2\to0\).

This is an \(S^3\) harmonic-level null test, not directly a sky-quadrupole-to-dipole ratio. The \(n=2\) eigenspace can project onto several observer-sky angular multipoles; an observational \(\ell_{\rm sky}=2\) prediction requires the explicit observer-shell projection and source weights.
